\section{Introduction}
\label{sec:introduction}
An agent that delegates work across organizations commits resources before it
knows how the work will unfold. It may reserve funds or authorize a tool call,
then discover a task requiring another provider. Each provider must be able to
determine which actions were authorized, which ran and which earned payment
as the plan changes.

Consider a buyer who funds an API survey and a follow-on review by an
intermediary agent. After the survey identifies a problem, the intermediary
hires a specialist from its own funds. The specialist completes its task and
earns payment, but the intermediary crashes before returning a result to the
buyer. Three facts must
survive the crash: the specialist's call has already run, its payment is owed,
and the buyer's task may still be incomplete. A replacement process must recover
those facts before deciding what to do next.

Plan changes introduce a related problem. A newly signed plan may contain both
the completed survey and the new specialist task. If the signature gives every
task a fresh execution right, the survey can run twice. If the new plan replaces
the evidence for the old one, a receiver may lose the terms under which the
survey ran. Valid signatures on each version do not, by themselves, establish
continuity between them.

We present Chio, a protocol and runtime kernel that preserves this continuity.
Its unit of composition is a \emph{work commitment}: signed records that bind
one selected invocation to a bounded allowance, the receiving owner's
authorization and, for paid work, its settlement terms. The records refer to
the same invocation by authenticated identifiers and digests. Each authority
checks the references relevant to its decision and retains the state needed
to prevent that decision from being spent twice.

The construction separates the identity of a task from the version of the plan
that contains it. An extension retains earlier tasks and allocations. A receiver
records the use of each execution right under a stable identity, so a signature
on a later plan cannot make that right unused again. Settlement reserves funds
for each agreement separately; once acceptance is recorded, a parent failure
cannot release a child's earned reserve. These rules permit incremental work
without a transaction that locks every participant at once.

Each receiving kernel admits work under its owner's capabilities, policies and
treaties. Foreign agreements supply evidence checked under locally selected
authority. We call this arrangement \emph{programmable sovereignty}: owners
express conditions for cooperation in verifiable agreements and enforce them
at their own resources.

We specify the records and transitions, establish a preservation result under
explicit trust assumptions, and implement the construction in Rust. Evaluation
follows the API example through graph growth and process loss, tests adversarial
substitutions, and compares the financial rules with transactional escrow and
ERC-8183.
